\documentclass[aspectratio=169, 11pt, xcolor={table}]{beamer}
% Shared style for the course slides, based on presentation_text/main.tex.
% Dark blue removed: every slide (cover included) uses the light background
% of the template's content slides, with teal + orange as accents.

\usepackage{amsmath}
\usepackage{amssymb}
\usepackage{booktabs}
\usepackage{colortbl}
\usepackage{tabularx}
\usepackage{tikz}
\usepackage[skins]{tcolorbox}
\usepackage{fontawesome5}
\usepackage{polyglossia}

\usetikzlibrary{positioning, calc, arrows.meta, shapes.geometric}

%------------------------------------------------
% Languages and fonts (compile from slides/ with XeLaTeX or tectonic)
%------------------------------------------------
\setdefaultlanguage[numerals=maghrib]{arabic}
\setotherlanguage{english}
\usefonttheme{serif}

\newfontfamily\arabicfont[
    Script=Arabic, Scale=0.98, Path=../fonts/, Extension=.ttf,
    UprightFont=*-Regular, BoldFont=*-Bold,
    ItalicFont=*-Italic, SlantedFont=*-Italic,
    BoldItalicFont=*-BoldItalic, BoldSlantedFont=*-BoldItalic]{Amiri}
\newfontfamily\arabicfontsf[
    Script=Arabic, Scale=0.98, Path=../fonts/, Extension=.ttf,
    UprightFont=*-Regular, BoldFont=*-Bold,
    ItalicFont=*-Italic, SlantedFont=*-Italic,
    BoldItalicFont=*-BoldItalic, BoldSlantedFont=*-BoldItalic]{Amiri}
% Heading font
\newfontfamily\khatt[Script=Arabic, Scale=1.05, Path=../fonts/,
    BoldFont=ae_AlBattar.ttf]{ae_AlBattar.ttf}

%------------------------------------------------
% Colors (template palette without the navy tones)
%------------------------------------------------
\definecolor{brandTeal}{HTML}{00B4D8}
\definecolor{brandOrange}{HTML}{FF7043}
\definecolor{brandLight}{HTML}{F8F9FA}
\definecolor{brandDark}{HTML}{263238}
\definecolor{brandMuted}{HTML}{6C757D}
\colorlet{tealDeep}{brandTeal!70!black}

\setbeamercolor{background canvas}{bg=brandLight}
\setbeamercolor{normal text}{fg=brandDark}
\setbeamercolor{structure}{fg=brandTeal}
\setbeamertemplate{navigation symbols}{}

% beamer.cls sets \raggedright globally; under RTL that centers short lines
% or pushes them left, so go back to justified text (flush right).
\rightskip=0pt\relax
% Lists call \raggedright too; make it a no-op for RTL slides.
\AtBeginDocument{\renewcommand{\raggedright}{\rightskip=0pt\relax}}

%------------------------------------------------
% Frame title, footer, lists
%------------------------------------------------
\setbeamertemplate{frametitle}{%
    \vspace{0.35cm}
    \noindent{\khatt\bfseries\Large\color{brandDark}\insertframetitle}
    \par
    \vskip-2pt
    \noindent{\color{brandTeal}\rule{\textwidth}{1.5pt}}
    \vspace{0.05cm}
}

% Footer: page number only
\setbeamertemplate{footline}{%
    \hbox to \paperwidth{\hfill
        {\color{brandMuted}\scriptsize\LR{\insertframenumber\,/\,\inserttotalframenumber}}%
        \hspace{0.6cm}}%
    \vspace{0.25cm}}

\setbeamertemplate{itemize item}{\color{brandTeal}\small\faChevronLeft}
\setbeamertemplate{itemize subitem}{\color{brandMuted}\small\faAngleLeft}
\setbeamertemplate{enumerate item}{\color{tealDeep}\bfseries\insertenumlabel)}

% Tables: every cell framed by a visible border
\newcolumntype{R}{>{\raggedleft\arraybackslash}X}
\newcolumntype{C}{>{\centering\arraybackslash}X}
\newcommand{\ras}[1]{\cellcolor{brandTeal!20}\textbf{#1}}
% Colored text inside a table cell (\leavevmode keeps it on the row baseline)
\newcommand{\cell}[2]{\leavevmode\textcolor{#1}{#2}}
\renewcommand{\arraystretch}{1.35}
\setlength{\arrayrulewidth}{0.6pt}
\arrayrulecolor{brandMuted!70}

%------------------------------------------------
% Boxes (the two boxes of the template's 2nd slide)
%------------------------------------------------
\tcbset{
    brandbox/.style={
        enhanced, arc=4pt, outer arc=4pt, boxrule=1.5pt,
        coltitle=white, coltext=brandDark, fonttitle=\khatt\bfseries,
        toptitle=0.8mm, bottomtitle=0.8mm,
        top=1.2mm, bottom=1.2mm, left=3mm, right=3mm,
        before skip=4pt, after skip=4pt}
}
% Teal box: definitions, objectives, summaries
\newtcolorbox{tbox}[2][]{brandbox, colback=brandTeal!10,
    colframe=brandTeal, title={#2}, #1}
% Orange box: warnings
\newtcolorbox{obox}[2][]{brandbox, colback=brandOrange!10,
    colframe=brandOrange, title={#2}, #1}
% Card (template's 3rd slide, in teal)
\newtcolorbox{card}[2][]{enhanced, colback=brandTeal!6, colframe=brandTeal,
    arc=6pt, boxrule=1.2pt, halign=center, halign title=center,
    coltitle=white, coltext=brandDark, fonttitle=\khatt\bfseries,
    top=1.5mm, bottom=1.5mm, left=2mm, right=2mm, title={#2}, #1}
% Prompt text written to the tool
\newtcolorbox{amr}[1][brandTeal]{enhanced, arc=1mm, boxrule=0pt,
    colback=white, overlay={\fill[#1] (frame.north east) rectangle ([xshift=-2.5pt]frame.south east);},
    left=6pt, right=8pt, top=2pt, bottom=2pt,
    before skip=3pt, after skip=3pt}

%------------------------------------------------
% TikZ (Arabic right-to-left inside nodes)
%------------------------------------------------
\newcommand{\nodecolor}{}
\tikzset{
    every node/.append style={
        execute at begin node={\nodecolor\everypar{\beginR}\beginR},
        execute at end node=\endR},
    lawn/.code={\def\nodecolor{\color{#1}}},
    khatwa/.style={draw=#1, fill=#1!12, very thick, rounded corners=5pt,
        minimum height=0.95cm, text width=1.9cm, align=center,
        font=\khatt\bfseries\small},
    khatwa/.default=brandTeal,
    sahm/.style={-{Stealth[length=2.5mm]}, thick, color=brandMuted},
    rang/.style={draw=#1, fill=#1!12},
    raqm/.style={circle, fill=#1, lawn=white, inner sep=1.5pt,
        font=\small\bfseries},
}

% Horizontal chain of steps, right to left:
% \flow[width]{{a}, {b\\c}, {d}}
\newcommand{\flow}[2][1.9cm]{%
    \begin{tikzpicture}[node distance=5mm]
        \foreach \t [count=\i] in {#2} {
            \ifnum\i=1
                \node[khatwa, text width=#1] (f\i) {\t};
            \else
                \pgfmathtruncatemacro\p{\i-1}
                \node[khatwa, text width=#1, left=of f\p] (f\i) {\t};
                \draw[sahm] (f\p) -- (f\i);
            \fi
        }
    \end{tikzpicture}}

%------------------------------------------------
% Cover (template's 1st slide on the light background)
%------------------------------------------------
\newcommand{\cover}[3]{%
\begin{frame}[plain]
    \begin{tikzpicture}[remember picture, overlay]
        \draw[brandTeal, line width=1.5pt, opacity=0.3]
            (current page.south east) ++(-1,1) circle (4cm);
        \draw[brandTeal, line width=0.8pt, opacity=0.25]
            (current page.north west) ++(2,2) circle (6cm);
        \draw[brandOrange, line width=0.5pt, opacity=0.25]
            (current page.north west) ++(2,2) circle (7cm);
        \node[anchor=north east] at ($(current page.north east)+(-1.0,-0.8)$) {
            \begin{tikzpicture}[scale=0.35]
                \fill[brandTeal] (0,0) -- (0.8,1.2) -- (1.6,0) -- cycle;
                \fill[brandOrange] (0.8,1.2) -- (1.6,2.4) -- (2.4,1.2) -- cycle;
            \end{tikzpicture}
        };
    \end{tikzpicture}

    \vspace{-0.05cm}
    {\hspace*{1.1cm}\khatt\bfseries\color{brandDark}
        المدرسة العليا للأساتذة ببوزريعة ــ ملحقة بومرداس\par}

    \vspace{1.2cm}
    \begin{minipage}{0.85\textwidth}
        {\color{tealDeep}\khatt\bfseries #1} \\[0.3cm]
        {\raggedleft\color{brandDark}\Huge\khatt\bfseries #2\par} \vspace{0.3cm}
        {\color{brandMuted}\large #3}
    \end{minipage}

    \vspace{1.2cm}
    \begin{minipage}{0.85\textwidth}
        \color{brandMuted}\small
        \begin{tabular}{@{}ll}
            \textbf{الأستاذ:} & د. بن عبد الرزاق يوسف \\
            \textbf{الفئة:}   & طلبة قسم اللغة العربية \\
        \end{tabular}
    \end{minipage}
\end{frame}}

\title{الأخلاقيات في الذكاء الاصطناعي}

\begin{document}

\cover{درس الفصل الرابع}
{الأخلاقيات في الذكاء الاصطناعي}
{المسؤولية ــ الشفافية ــ الإفصاح ــ الأمانة العلمية}

\begin{frame}{أهداف الدرس}
    \begin{tbox}{\faBullseye\quad أهداف الدرس}
        في نهاية هذا الدرس، تستطيع أن:
        \begin{itemize}
            \item تعرّف المسؤولية والشفافية والإفصاح.
            \item تميّز بين الاستعمال المقبول والاستعمال المرفوض.
            \item تعرّف الانتحال العلمي والانتحال بالذكاء الاصطناعي.
            \item تكتب جملة إفصاح بسيطة في عملك.
        \end{itemize}
    \end{tbox}
\end{frame}

\begin{frame}{لماذا نتحدث عن الأخلاقيات؟}
    الذكاء الاصطناعي أداة قوية وسريعة، لكن الأداة لا تقرر وحدها كيف نستعملها.

    \vspace{3mm}
    \begin{tbox}{\faUserCheck\quad الإنسان هو الذي يختار}
        يستعمل الأداة بأمانة، أو يستعملها للغش.\\[3pt]
        وأستاذ المستقبل \textbf{قدوة} لتلاميذه، فسلوكه مع هذه الأدوات درس في حد ذاته.
    \end{tbox}
\end{frame}

\begin{frame}{ثلاثة مبادئ أساسية}
    \begin{center}
    \begin{tikzpicture}[
        amud/.style={draw=brandTeal, fill=brandTeal!12, very thick,
            minimum width=3.4cm, minimum height=2.3cm,
            text width=3.1cm, align=center}]
        \fill[brandOrange] (-6,1.9) -- (6,1.9) -- (0,3.1) -- cycle;
        \node[lawn=white, font=\khatt\bfseries] at (0,2.3) {الاستعمال الأخلاقي};
        \node[amud] at (4,0.55)
            {{\khatt\bfseries المسؤولية}\\[3pt]\small أنا مسؤول عمّا أقدّمه};
        \node[amud] at (0,0.55)
            {{\khatt\bfseries الشفافية}\\[3pt]\small أوضّح كيف استعملت الأداة};
        \node[amud] at (-4,0.55)
            {{\khatt\bfseries الإفصاح}\\[3pt]\small أصرّح باستعمالها};
        \fill[brandMuted!70] (-6,-0.75) rectangle (6,-1.1);
    \end{tikzpicture}
    \end{center}
\end{frame}

\begin{frame}{تعريف المبادئ}
    \begin{tbox}{\faLightbulb\quad تعريف}
        \textbf{المسؤولية:} المستخدم مسؤول عن استعمال الأداة، وعن كل معلومة يعتمد عليها.\\[2pt]
        \textbf{الشفافية:} توضيح طريقة استعمال الذكاء الاصطناعي عند الحاجة.\\[2pt]
        \textbf{الإفصاح:} التصريح باستعمال الأداة عندما يكون ذلك مطلوبًا أو مناسبًا.
    \end{tbox}
    \begin{obox}{\faExclamationTriangle\quad انتبه!}
        الاستعانة بالذكاء الاصطناعي \textbf{لا تنقل المسؤولية} إلى الأداة.
        إذا قدّمت معلومة خاطئة، فأنت المسؤول عنها، لا الأداة.
    \end{obox}
\end{frame}

\begin{frame}{الانتحال العلمي}
    \begin{tbox}{\faLightbulb\quad تعريف}
        \textbf{الانتحال العلمي:} أن ينسب الشخص إلى نفسه عملًا أو أفكارًا ليست له،
        دون ذكر مصدرها.\\[3pt]
        \textbf{الانتحال بالذكاء الاصطناعي} \textenglish{(AI-Plagiarism)}:
        أن ينسخ الطالب نصًا أنتجته الأداة، ويقدمه على أنه عمله الشخصي.
    \end{tbox}
\end{frame}

\begin{frame}{حالة للتحليل}
    \begin{obox}{\faExclamationTriangle\quad حالة}
        نسخ طالب بحثًا كاملًا من أداة ذكاء اصطناعي، ثم كتب اسمه عليه وقدّمه للأستاذ.
    \end{obox}
    هذا \textbf{انتحال}، لأن:
    \begin{itemize}
        \item العمل ليس عمله.
        \item لم يتحقق من صحته.
        \item لم يصرّح بمصدره.
    \end{itemize}
\end{frame}

\begin{frame}{مقياس الاستعمال: من المقبول إلى المرفوض}
    \vspace*{-2mm}
    \begin{center}
    \begin{tikzpicture}
        \shade[left color=brandOrange, right color=brandTeal, rounded corners=3pt]
            (-6,0) rectangle (6,0.4);
        \node[font=\khatt\bfseries, lawn=tealDeep] at (5,0.75)  {مقبول};
        \node[font=\khatt\bfseries, lawn=brandMuted] at (0,0.75) {مقبول بشروط};
        \node[font=\khatt\bfseries, lawn=brandOrange] at (-5,0.75) {مرفوض};
    \end{tikzpicture}
    \end{center}
    \small\renewcommand{\arraystretch}{1.15}
    \begin{tabularx}{\textwidth}{|R|>{\centering\arraybackslash}p{3.2cm}|}
        \hline
        \ras{الاستعمال} & \ras{الموقف الأخلاقي} \\ \hline
        استعمال الأداة لفهم قاعدة نحوية صعبة.
            & \cell{tealDeep}{\textbf{مقبول}} \\ \hline
        طلب أفكار من الأداة، ثم كتابة الإجابة بنفسي.
            & \cell{tealDeep}{\textbf{مقبول}} \\ \hline
        إنشاء مسودة بالأداة، ثم مراجعتها والتحقق منها والإفصاح.
            & \cell{brandMuted}{\textbf{مقبول بشروط}} \\ \hline
        الاعتماد على معلومات الأداة دون التحقق منها.
            & \cell{brandOrange}{\textbf{غير مسؤول}} \\ \hline
        نسخ إجابة كاملة وتقديمها باسمي.
            & \cell{brandOrange}{\textbf{انتحال مرفوض}} \\ \hline
    \end{tabularx}
\end{frame}

\begin{frame}{كيف أكتب جملة إفصاح؟}
    جملة الإفصاح قصيرة وصادقة. نضعها في آخر البحث أو في مقدمته.
    \begin{amr}[brandTeal]
        \textbf{نموذج:} استعملتُ أداة \textenglish{ChatGPT} لتنظيم أفكار الفصل الثاني
        واقتراح خطة أولية. راجعتُ كل المعلومات وتحققت منها في المصادر المذكورة،
        والصياغة النهائية من عملي.
    \end{amr}
    \begin{tbox}{\faChalkboardTeacher\quad لتلاميذك}
        علّمهم هذه القاعدة البسيطة:
        «الأداة تساعدني على التفكير، ولا تفكر بدلًا مني».
    \end{tbox}
\end{frame}

\begin{frame}{قواعد ذهبية للاستعمال الأخلاقي}
    \begin{center}
    \begin{tikzpicture}[
        qaida/.style={draw=brandTeal, fill=brandTeal!10, thick, rounded corners=4pt,
            text width=10cm, align=right, minimum height=0.85cm, inner sep=5pt}]
        \node[qaida] (q1) at (0,0)     {أتحقق من كل معلومة مهمة قبل استعمالها.};
        \node[qaida] (q2) at (0,-1.05) {أكتب العمل النهائي بأسلوبي وأفكاري.};
        \node[qaida] (q3) at (0,-2.1)  {أصرّح باستعمال الأداة عندما يكون ذلك مطلوبًا.};
        \node[qaida] (q4) at (0,-3.15) {أحترم قوانين المؤسسة وتعليمات الأستاذ.};
        \foreach \i in {1,...,4}
            \node[raqm=brandOrange, minimum size=0.6cm] at ($(q\i.east)+(0.5,0)$) {\i};
    \end{tikzpicture}
    \end{center}
\end{frame}

\begin{frame}{خلاصة الدرس}
    \begin{tbox}{\faCheckCircle\quad خلاصة الدرس}
        \begin{itemize}
            \item المبادئ الثلاثة: \textbf{المسؤولية}، \textbf{الشفافية}، \textbf{الإفصاح}.
            \item الأداة لا تتحمل المسؤولية؛ \textbf{المستخدم} هو المسؤول.
            \item نسخ عمل الأداة وتقديمه باسمي \textbf{انتحال}.
            \item الاستعمال الجيد: أفهم، أتحقق، أكتب بنفسي، وأصرّح.
        \end{itemize}
    \end{tbox}
\end{frame}

\end{document}
